\documentclass[12pt]{article}
\usepackage[T1]{fontenc}
\usepackage[utf8]{inputenc}
\usepackage{lmodern}
\usepackage{textcomp}
\usepackage{enumitem}
\usepackage{geometry}
\geometry{margin=1in}
\begin{document}
\begin{center}
Answer Key - History - Undergraduate Level Assessment 11 \
\small (Answers are ordered exactly as in the question paper.)
\end{center}

\section*{Multiple Choice}
\begin{enumerate}[label=\arabic*.]
\item \textbf{Answer:} B
\\ \textit{Options:} A) 1.8 million; 1.2 million; B) 2.5 million; 1.8 million; C) 1.4 million; 800,000; D) 4.5 million; 3.2 million; E) 2 million; 1.5 million
\\ \textit{Note:} The correct answer reflects the established timeline in human evolution, where Homo habilis is believed to have emerged approximately 2.5 million years ago, followed by Homo erectus around 1.8 million years ago, based on fossil evidence and archaeological findings.
\item \textbf{Answer:} D
\\ \textit{Options:} A) severe drought and a surge in warfare; B) increased rainfall and population migration to the Maya ceremonial centers; C) severe drought and increased storage of surplus food from agriculture; D) severe drought and a decline in construction at Maya ceremonial centers; E) increased warfare and the abandonment of the entire region
\\ \textit{Note:} The correct answer is supported by the evidence that periods of severe drought, indicated by variations in oxygen isotopes in stalagmites, led to decreased agricultural productivity, which in turn caused a decline in the construction and maintenance of Maya ceremonial centers.
\item \textbf{Answer:} A
\\ \textit{Options:} A) use of a calendar system based on celestial bodies; B) practice of mummification for the dead; C) ritual consumption of maize beer; D) construction of underground dwellings; E) ritual human sacrifice
\\ \textit{Note:} The Hohokam adopted a calendar system based on celestial bodies from Mesoamerican culture, as evidenced by their agricultural practices and ceremonial activities that aligned with astronomical events.
\item \textbf{Answer:} A
\\ \textit{Options:} A) Sunda; Sahul; B) Sunda; Wallacea; C) Gondwana; Beringia; D) Sunda; Gondwana; E) Beringia; Sunda
\\ \textit{Note:} The correct answer is based on the understanding that during the Pleistocene era, the landmasses of Java, Sumatra, Bali, and Borneo were part of a larger continental shelf known as Sundaland, while the land connecting Australia, New Guinea, and Tasmania was referred to as Sahul, highlighting the geographical and geological distinctions of these regions during that period.
\item \textbf{Answer:} A
\\ \textit{Options:} A) the discovery of an artifact in a museum or private collection; B) an archaeological site found underwater; C) a site where artifacts have been intentionally buried for preservation; D) all of the above; E) a site where artifacts have been moved and rearranged by humans or natural forces
\\ \textit{Note:} The discovery of an artifact in a museum or private collection indicates a primary context because it provides direct evidence of the artifact's original use and significance, allowing historians to analyze it in relation to its cultural and historical background.
\end{enumerate}

\section*{True / False}
\begin{enumerate}[label=\arabic*.]
\setcounter{enumi}{5}
\item \textbf{Answer:} True
\\ \textit{Options:} A) True; B) False
\\ \textit{Note:} According to the passage: Two of the expeditions were successful, with Fort Duquesne and Louisbourg
\item \textbf{Answer:} False
\\ \textit{Options:} A) True; B) False
\\ \textit{Note:} The correct answer is: May 3
\item \textbf{Answer:} False
\\ \textit{Options:} A) True; B) False
\\ \textit{Note:} The correct answer is: Protestant Reformation
\item \textbf{Answer:} False
\\ \textit{Options:} A) True; B) False
\\ \textit{Note:} The correct answer is: Adam Johann von Krusenstern
\item \textbf{Answer:} False
\\ \textit{Options:} A) True; B) False
\\ \textit{Note:} The correct answer is: 999
\end{enumerate}

\section*{Long Form Response}
\begin{enumerate}[label=\arabic*.]
\setcounter{enumi}{10}
\item \textbf{Answer:} 1990
\\ \textit{Note:} Expected answer: 1990
\item \textbf{Answer:} proper names of the kings of Kish
\\ \textit{Note:} Expected answer: proper names of the kings of Kish
\end{enumerate}

\vfill
\noindent\textit{End of Answer Key}
\end{document}